% This file is part of proScientia.ltx
% (c) 2022 Karsten Reincke (https://github.com/kreincke/proScientia.ltx)
% It is distributed under the terms of the creative commons license
% CC-BY-4.0 (= https://creativecommons.org/licenses/by/4.0/)

\documentclass[
  DIV=calc,
  BCOR=5mm,
  11pt,
  headings=small,
  oneside,
  abstract=true,
  toc=bib,
  english,ngerman]{scrbook}

%%% (1) general configurations %%%
\usepackage[utf8]{inputenc}
\usepackage{a4}
\usepackage[ngerman]{babel}

%language specific quoting signs
\usepackage[
  style=german,
  autostyle=true,
]{csquotes}

% language specific hyphenation
% This file is part of proScientia.ltx
% (c) 2022 Karsten Reincke (https://github.com/kreincke/proScientia.ltx)
% It is distributed under the terms of the creative commons license
% CC-BY-4.0 (= https://creativecommons.org/licenses/by/4.0/)

\hyphenation{ pro-scien-tia there-fo-re}

\babelhyphenation{ pro-scien-tia wor-tmi-tfalsch-entr-ennu-ngen}


%%% (3) layout page configuration %%%

% select the visible parts of a page
% S.31: { plain|empty|headings|myheadings }
\pagestyle{headings}

% select the wished style of page-numbering
% S.32: { arabic,roman,Roman,alph,Alph }
\pagenumbering{arabic}
\setcounter{page}{1}

% select the wished distances using the general setlength order:
% S.34 { baselineskip| parskip | parindent }
% - general no indent for paragraphs
\setlength{\parindent}{0pt}
\setlength{\parskip}{1.2ex plus 0.2ex minus 0.2ex}

% Hyperlinks
\usepackage{hyperref}
\hypersetup{bookmarks=true,breaklinks=true,colorlinks=true,citecolor=blue,draft=false}

\usepackage[T1]{fontenc}
\usepackage{wesu}
\usepackage{suetterl} % provided by fundus-suetterlin of TexLive

\usepackage{fontawesome6}

\begin{document}

%% use all entries of the bliography
\nocite{*}


{\Large \wesu 

Ein Professor erzählte einmal wie er vor einigen Jahren vielleicht waren es sieben oder acht 
eine Einführungs:vorlesung gehalten habe Wie er es immer täte habe er auch damals mit dem 
Thema Logik begonnen Daher habe er zunächst erklärt was eine Aussage sei nämlich ein Satz 
oder ein Text dessen Inhalt entweder wahr oder falsch sei Als Beispiel habe er 
ohne gro{\ss} nachzudenken den Satz genannt Karl ist krank 

Kurz nachdem er so der Professor weiter das Beispiel ausgesprochen habe sei ihm jedoch 
siedendhei{\ss} eingefallen dass er unbedingt eine lebende Person benötigt habe die Karl hie{\ss} 
sonst hätte sich sein Beispiel auf niemanden bezogen und der Satz wäre weder wahr noch falsch 
gewesen So etwas meinte der Professor sei ihm zuvor noch nie passiert und die Situation sei 
ihm besonders peinlich gewesen da es sich um einen Anfängerfehler handele 

Um seinen unüberlegten Fehler schnell wieder auszubügeln habe er dann in den Hörsaal hinein 
gefragt ob jemand anwesend sei der den Namen Karl trage Die Studenten hätten zunächst irritiert 
geschaut weil ihnen nicht klar gewesen sei warum dies von Belang gewesen wäre Sekundenlang 
habe Stille geherrscht bis sich aus einer der letzten Reihen jemand mit der 
Bemerkung zu Wort habe Der ist krank

}

\vspace{0.5cm}

\footnotesize{
Entnommen aus \href{https://www.diktat-truhe.de/}{https://www.diktat-truhe.de/}. In die indirekte Rede 'übersetzt' v. K.R.
}
\end{document}
